%%%%%%%%%%%%%%%
% CV of Gerardo Cicalese
% Clean academic version for PhD and research use
\documentclass[a4paper,skipsamekey,11pt,english]{curve}

\PassOptionsToPackage{style=ieee,sorting=ydnt,uniquename=init,defernumbers=true}{biblatex}
\usepackage{settings}
\DefineBibliographyStrings{english}{url={\textsc{url}}}

\mynames{Cicalese/Gerardo}

\ifxetexorluatex
  \usepackage{fontspec}
  \usepackage[p,osf,swashQ]{cochineal}
  \usepackage[medium,bold]{cabin}
  \usepackage[varqu,varl,scale=0.9]{zi4}
\else
  \usepackage[T1]{fontenc}
  \usepackage[p,osf,swashQ]{cochineal}
  \usepackage{cabin}
  \usepackage[varqu,varl,scale=0.9]{zi4}
\fi

\includecomment{fullonly}

%%%%%%%%%%%%%%%%%%%%%%%%%%%%%%%%%%%%%%

\leftheader{%
  {\LARGE\bfseries\sffamily Gerardo Cicalese}

  \makefield{\faUniversity}{Politecnico di Milano}
  \makefield{\faGraduationCap}{Ph.D. Student in Mathematical Models and Methods in Engineering}

  % Optional contact info:
  % \makefield{\faEnvelope[regular]}{\href{mailto:gerardo.cicalese@example.com}{\texttt{gerardo.cicalese@example.com}}}
  % \makefield{\faLinkedin}{\href{https://www.linkedin.com/in/gerardocicalese}{\texttt{linkedin.com/in/gerardocicalese}}}
}

\rightheader{}

\begin{fullonly}
\photo[r]{photo}
\photoscale{0.13}
\end{fullonly}

\title{Curriculum Vitae}

\begin{document}
\makeheaders[c]

%%% --- EDUCATION ---------------------------------------------------------
\begin{rubric}{Education}
\entry*[2023--present]%
	\textbf{Ph.D. in Mathematical Models and Methods in Engineering}, Politecnico di Milano, Italy.
	\par Thesis: \emph{Hybrid numerical methods for the design of acoustically comfortable environments.}

\entry*[2021--2023]%
	\textbf{M.Sc. in Music and Acoustic Engineering}, Politecnico di Milano, Italy.
	\par Thesis: \emph{Efficient techniques for accurate sound propagation modeling in virtual acoustics.}
	\par Final mark: 110L/110.

\entry*[2018--2021]%
	\textbf{B.Sc. in Computer Engineering}, University of Salerno, Italy.
	\par Thesis: \emph{Analisi e test di Jukebox, una piattaforma per la generazione di pattern musicali.}
	\par Final mark: 110L/110.

\entry*[2013--2018]%
	\textbf{Scientific High School Diploma}, Liceo Scientifico ``Nicola Sensale'', Nocera Inferiore, Italy.
	\par Final mark: 98/100.
\end{rubric}

%%% --- RESEARCH INTERESTS ------------------------------------------------
\begin{rubric}{Research Interests}
\entry*[2023--present]%
	Numerical methods for wave propagation and acoustics
	\par Domain decomposition and waveform relaxation methods  
	\par Parallel computing for scientific simulations
	\par Virtual/augmented reality and auralization
	\par Machine learning for acoustic field prediction
\end{rubric}

%%% --- PUBLICATIONS ------------------------------------------------------
\begin{rubric}{Publications}
\entry*[2024]%
	\textbf{Addressing Atmospheric Absorption in Adaptive Rectangular Decomposition.}
	\par \textit{The Journal of the Acoustical Society of America}, 156(4):2328--2339.
	\par Gerardo Cicalese, Gabriele Ciaramella, Ilario Mazzieri.

\entry*[In Prep.]%
	\textbf{Optimized Schwarz Waveform Relaxation for the Damped Wave Equation.}
	\par Gerardo Cicalese, Gabriele Ciaramella, Ilario Mazzieri, Martin J. Gander.
\end{rubric}

%%% --- CONFERENCES -------------------------------------------------------
\begin{rubric}{Conferences and Presentations}
\entry*[2025]%
	\textbf{Swiss Numerical Analysis Day 2025}, Basel
	\par Poster: \emph{Enhancing Adaptive Rectangular Decomposition with Atmospheric Absorption Modeling}

\entry*[2025]%
	\textbf{Domain Decomposition 29 (DD29)}, Milan
	\par Poster: \emph{Enhancing Adaptive Rectangular Decomposition with Atmospheric Absorption Modeling}
	\par Talk: \emph{Optimized Red--Black Waveform Relaxation for the Damped Wave Equation}

\entry*[2025]%
	\textbf{SIMAI 2025}, Trieste
	\par Talk: \emph{Optimized Red--Black Waveform Relaxation for the Damped Wave Equation}
\end{rubric}

%%% --- RESEARCH STAYS ----------------------------------------------------
\begin{rubric}{Research Stays}
\entry*[2025]%
	\textbf{Université de Genève}, Switzerland
	\par Research visits supervised by Prof. Martin J. Gander
	\par March--May 2025; September--November 2025
\end{rubric}

%%% --- TEACHING ----------------------------------------------------------
\begin{rubric}{Teaching Experience}
\entry*[2023--2025]%
	\textbf{Teaching Assistant}, \emph{Curve e superfici per il design}
	\par Prof. Gabriele Ciaramella, Politecnico di Milano

\entry*[2023--2024]%
	\textbf{Teaching Assistant}, \emph{Calcolo numerico}
	\par Prof. Carlo De Falco, Politecnico di Milano
\end{rubric}

%%% --- SKILLS ------------------------------------------------------------
\begin{rubric}{Skills}
\entry*[Programming]%
	MATLAB, Python, C/C++, \LaTeX

\entry*[Scientific Tools]%
	COMSOL Multiphysics, Intel VTune Profiler, OpenMP

\entry*[Numerical Methods]%
	Finite Difference Time Domain (FDTD), Finite Element Method (FEM), Domain Decomposition, Waveform Relaxation

\entry*[Languages]%
	Italian (native), English (fluent, C1 level)
\end{rubric}

%%% --- END ---------------------------------------------------------------
\end{document}